\section{Project Overview}

\subsection{Engineering objective}
The engineering problem is to tune a practical finite-radius wire dipole so that its
input match occurs near \SI{1}{GHz}, while retaining the expected broadside radiation
pattern of a resonant half-wave dipole. The project is intentionally framed as a compact
verification case: analytical theory supplies a canonical baseline, CST provides the
full-wave response, and the comparison identifies both agreement and missing evidence.

\subsection{Design requirements and configuration}
\begin{table}[H]
\centering
\caption{Design requirements and documented model parameters.}
\label{tab:requirements}
\begin{tabularx}{\textwidth}{>{\bfseries}p{0.31\textwidth}X}
\toprule
Item & Documented value or requirement \\
\midrule
Target frequency & \SI{1.000}{GHz} \\
Free-space wavelength used for design & \SI{300}{mm} \\
Initial theoretical length & \SI{150}{mm} \\
Final tuned total length & \SI{135.6}{mm} \\
Conductor model & PEC cylinder, \SI{5}{mm} diameter \\
Center feed gap & \SI{20}{mm} \\
Port reference impedance & \SI{73}{\ohm} \\
Frequency sweep & \SIrange{0.5}{2.0}{GHz} \\
Far-field monitors & \SI{0.9}{GHz} and \SI{1.0}{GHz} \\
Primary outputs & $S_{11}$, 2D pattern, 3D directivity pattern \\
\bottomrule
\end{tabularx}
\end{table}

\subsection{Key engineering outcome}
The final \SI{135.6}{mm} model produced the exact CST marker result
$f_{\mathrm{res}}=\SI{0.99648}{GHz}$ and
$S_{11,\min}=-\SI{47.01186}{dB}$. Its displayed directivity at \SI{1}{GHz} was
\SI{2.149}{dBi}, essentially equal to the analytical half-wave value. These results
support a tuned resonant dipole interpretation under the documented \SI{73}{\ohm}
reference condition.

\subsection{Reporting decision}
The marker-based final run is used as the quantitative design result. The earlier
\SI{150}{mm} and \SI{125}{mm} cases are treated as directional tuning evidence because
their numerical minima are available only from plotted curves, not raw exported data.
